%==============================================================================
% Chapter 7 — Computer Vision & Scientific ML
% Curriculum: Phase 2, Chapter 7 of docs/AI_ML_Master_Checklist.md
% Companion code: projects/07_cv_sciml/
%
% Section skeleton only. Figures for this chapter go in theory/figures/.
%==============================================================================
\chapter{Computer Vision \& Scientific ML}
\label{ch:cv-sciml}

Two applications of \cref{ch:deep-learning} that look unrelated and share a
central idea: both get their power from building a known structure of the
problem directly into the model rather than hoping the network discovers it.\projectref{projects/07\_cv\_sciml/}

In computer vision that structure is spatial --- translation equivariance and
locality, encoded as weight sharing in a convolution. In scientific machine
learning it is physical --- a differential equation the solution must satisfy,
encoded as an extra term in the loss. Physics-informed networks are the clearer
illustration: the PDE residual is computed by automatic differentiation of the
network with respect to its \emph{inputs}, which is the same machinery as
backpropagation pointed in an unfamiliar direction.

%==============================================================================
\section{Convolutional Neural Networks}
\label{sec:cnn}

\subsection{The Discrete Convolution}

\subsection{Strides, Padding and Output Dimensions}

\subsection{Pooling Layers}

\subsection{Receptive Fields}

%==============================================================================
\section{Classic Architectures}
\label{sec:cv-architectures}

\subsection{VGG and the Depth Hypothesis}

\subsection{ResNet and the Skip Connection}

%==============================================================================
\section{Object Detection and Segmentation}
\label{sec:detection-segmentation}

\subsection{Intersection over Union}

\subsection{YOLO}

\subsection{Mask R-CNN}

%==============================================================================
\section{Modern Computer Vision}
\label{sec:modern-cv}

\subsection{Vision Transformers}

\subsection{Patch Embedding}

\subsection{Transfer Learning and Fine-Tuning}

%==============================================================================
\section{Physics-Informed Neural Networks}
\label{sec:pinn}

\subsection{Embedding a PDE in the Loss Function}

\subsection{Collocation Points and Boundary Conditions}

\subsection{Balancing the Data and Physics Terms}

%==============================================================================
\section{Neural Operators}
\label{sec:neural-operators}

\subsection{Learning Mappings Between Function Spaces}

\subsection{The Fourier Neural Operator}

\subsection{Solving PDEs in the Frequency Domain}
